\documentclass[10pt]{article}

\usepackage[a4paper, total={6in, 9.5in}]{geometry}
\usepackage[dvipsnames,svgnames]{xcolor}
\usepackage[shortlabels]{enumitem}
\usepackage[framemethod=TikZ]{mdframed}
\usepackage{amsmath,amssymb,amsthm}
\usepackage[colorlinks]{hyperref}
%\usepackage{draftwatermark}
%\SetWatermarkScale{7}
%\SetWatermarkColor[gray]{0.9}

\hypersetup{
	colorlinks=true,
	linkcolor=blue,
	filecolor=blue,      
	urlcolor=blue,
	pdftitle={Aditya Pahuja's CV}
}

\usepackage[textsize=scriptsize]{todonotes}
\usepackage{nopageno}

\setlist[itemize]{label={$\circ$},itemsep=0pt,topsep=2pt}

\renewcommand{\tilde}{\widetilde}
\renewcommand{\hat}{\widehat}
\renewcommand{\setminus}{\!-\!}
\providecommand{\ol}{\overline}
\providecommand{\eps}{\varepsilon}
\providecommand{\half}{\frac{1}{2}}
\providecommand{\dang}{\measuredangle} %% Directed angle
\providecommand{\CC}{\mathbb C}
\providecommand{\FF}{\mathbb F}
\providecommand{\HH}{\mathbb H}
\providecommand{\NN}{\mathbb N}
\providecommand{\QQ}{\mathbb Q}
\providecommand{\RR}{\mathbb R}
\providecommand{\ZZ}{\mathbb Z}
\providecommand{\RP}{\mathbb RP}
\providecommand{\CP}{\mathbb CP}
\providecommand{\dg}{^\circ}
\providecommand{\ind}[1]{\mathbf 1_{#1}}
\providecommand{\ii}{\item}
\providecommand{\alert}[1]{{\sffamily\textbf{\textcolor{magenta}{#1}}}}
\providecommand{\opname}{\operatorname}
\providecommand{\li}{\opname{li}}
\providecommand{\ls}{\opname{ls}}
\providecommand{\supp}{\opname{supp}}
\providecommand{\inv}{^{-1}}
\providecommand{\setdiff}{\smallsetminus}
\providecommand{\mb}[1]{\mathbf{#1}}
\providecommand{\abs}[1]{\left\lvert{#1}\right\rvert}
\providecommand{\norm}[1]{\left\|{#1}\right\|}
\providecommand{\dif}{\;d}
\providecommand{\id}{\operatorname{id}}

\setlength{\parindent}{0pt}

\title{Aditya Pahuja\vspace{-1em}}
%\author{\normalsize aditya.pahuja@stonybrook.edu
%$\spadesuit$ \url{https://hungryforpi.github.io/}}
\date{\normalsize Last updated \today}

\makeatletter
\renewcommand{\@maketitle}{
    \newpage
    \begin{center}
    \vspace*{0em} % This removes the initial vertical space
    {\LARGE \@title \par}
    \vspace{0.5em} % Space between title and author
    {\large
      \lineskip .5em
      \begin{tabular}[t]{c}\@author\end{tabular}\par}
    \vspace{0.5em} % Space between author and date
    {\large \@date\par}
    \end{center}
    \par
}
\makeatother

\begin{document}
\maketitle

\vspace{-3mm}\section*{Education}
\textbf{Stony Brook University, B.S.\ in Mathematics}
\hfill{}August 2024 -- May 2027

GPA: 3.96/4.00; Major GPA: 4.00/4.00
\vspace{1mm}

Relevant coursework:
\begin{itemize}
    \small
    % don't include too many undergrad ones
    \ii Undergraduate-level:
    multivariable calculus and linear algebra;
    real analysis;
    analysis in several variables;
    differential geometry of curves and surfaces;
    abstract algebra II
    % reorder these? e.g. alphabetically
    \ii Graduate-level:
    foundations of topology;
    algebra I;
    geometry/topology II;
    algebra II;
    complex analysis
\end{itemize}

\textbf{Stuyvesant High School}
\hfill{}September 2020 -- June 2024

\vspace{-3mm}\section*{Research Experience}
%\textbf{REU at University of Chicago}
%\hfill{}June 2026 -- August 2026
%\begin{itemize}
%    \ii[*] But nobody came.
%\end{itemize}

\textbf{Directed Reading Program Participant}
\hfill{}January 2026 -- May 2026
\begin{itemize}
    \ii Read about differential topology independently
    and discussed it weekly with a graduate student mentor.
    \ii Gave a short presentation at the end of the semester
    about vector fields on the sphere and wrote an expository paper
    on constructing characteristic classes via obstruction theory.
    % TODO do and link
\end{itemize}

\textbf{REU on machine learning in computational algebra}
\hfill{}June 2025 -- December 2025
\begin{itemize}
    \ii Used SageMath and PyTorch to generate data for and train neural networks
    to predict runtime for Gröbner basis computation of a given polynomial ideal
    and term ordering.
\end{itemize}

\vspace{-3mm}\section*{Work Experience}
\textbf{NYC Interscholastic Math League}

\emph{Coordinator and Systems Administrator}
\hfill{} June 2025 -- present
\begin{itemize}
    \ii Supervised and guided the Soph-Frosh and Senior A
    contest writing teams during the fall and spring semesters
    to ensure quality and appropriate difficulty of
    three Soph-Frosh and five Senior A contests
    throughout the school year.
    \ii Restructured the \LaTeX{} files used to generate contests
    in order to organize them more modularly,
    making collaborative editing more convenient
    and reducing the chance of human error.
\end{itemize}

\emph{Senior A Problem Writer}
\hfill{}June 2024 -- May 2025

\begin{itemize}
    \ii Collaborated to write 10 high-quality math contests for
    high school seniors throughout NYC.
\end{itemize}

\textbf{Teaching Assistant, MAT 125}
\hfill{} August 2025 -- December 2025
\begin{itemize}
    \ii Planned and taught a weekly 50-minute recitation section
    for 28 students in an introductory calculus course
    to supplement the lectures.
    \ii Helped students from various math courses,
    including calculus, intro to proofs, and real analysis,
    in the Math Learning Center for three hours a week.
\end{itemize}

\vspace{-3mm}\section*{Volunteering}
\textbf{New York City Math Team}
\hfill{}September 2024 -- present

\begin{itemize}
    \ii Supervised Saturday math team practices during the spring
    and wrote a mock contest to provide the team with practice materials
    for the ARML math competition.
    \ii Represented NYC as a coach at NYSML 2025 and 2026,
    where I respectively proctored a team and helped grade the power round.
    %\ii Occasional guest lecturer for Stuyvesant High School’s senior math team.
\end{itemize}

\vspace{-3mm}\section*{Honors and Awards}
\textbf{Kuga-Sah Memorial Award}
\hfill{}received May 2025, May 2026
\begin{itemize}
    \ii Award given to two students in each grade level
    at Stony Brook University for excellence in mathematics.
\end{itemize}
\textbf{Honorable Mention, Putnam Mathematics Competition}
\hfill{}received February 2025
\begin{itemize}
    \ii Scored 48/120 on the 2024 Putnam math contest, achieving rank 100
    among 3988 participants.
\end{itemize}

%\vspace{-3mm}\section*{Skills and Miscellanea}
%% TODO organize
%Languages: English, French, Hindi % TODO include level of skillz
%
%Programming languages: Python, HTML, \LaTeX{}
%
%% google sheets?










\end{document}
